\chapter{Πλήρες Σύνολο Χαρτών Συσταδοποίησης}
\label{app:amb:cluster-maps}

Το παράρτημα παρουσιάζει τα selected cluster maps για όλους τους
συνδυασμούς order sweep, περιοχής ελέγχου και μεθόδου συσταδοποίησης.
Οι χάρτες αντιστοιχούν στις επιλεγμένες ρυθμίσεις του Κεφαλαίου~3.

\newcommand{\ambientappendixmap}[5]{%
\begin{figure}[p]
    \centering
    \includegraphics[width=0.94\textwidth,height=0.80\textheight,keepaspectratio]{\ambientfigurepath Ambient_Mag0.1_T600s_dt10ms_seed1997/#1/clustering/by_control_area/#2/#3/pdf/#4}
    \caption{Selected cluster map: #5 (#1).}
\end{figure}
}

\newcommand{\ambientappendixcase}[3]{%
\section{#3 \textendash{} #1}
\ambientappendixmap{#1}{#2}{kmeans}{k-means_selected_cluster_map.pdf}{\textit{k-Means}}
\ambientappendixmap{#1}{#2}{kmedoids}{k-medoids_selected_cluster_map.pdf}{\textit{k-Medoids}}
\ambientappendixmap{#1}{#2}{dbscan}{dbscan_selected_cluster_map.pdf}{DBSCAN}
\ambientappendixmap{#1}{#2}{optics}{optics_selected_cluster_map.pdf}{OPTICS}
\ambientappendixmap{#1}{#2}{hdbscan}{hdbscan_selected_cluster_map.pdf}{HDBSCAN}
\ambientappendixmap{#1}{#2}{gmm}{gmm_selected_cluster_map.pdf}{GMM}
\ambientappendixmap{#1}{#2}{agglomerative}{agglomerative_selected_cluster_map.pdf}{Agglomerative Hierarchical Clustering}
\clearpage
}

\ambientappendixcase{orders1}{area_1}{Περιοχή 1}
\ambientappendixcase{orders1}{area_2}{Περιοχή 2}
\ambientappendixcase{orders1}{area_3}{Περιοχή 3}
\ambientappendixcase{orders2}{area_1}{Περιοχή 1}
\ambientappendixcase{orders2}{area_2}{Περιοχή 2}
\ambientappendixcase{orders2}{area_3}{Περιοχή 3}
